\documentclass[11pt,a4paper]{article}

% Packages
\usepackage[utf8]{inputenc}
\usepackage[T1]{fontenc}
\usepackage{amsmath,amsfonts,amssymb}
\usepackage{graphicx}
\usepackage{booktabs}
\usepackage{enumitem}
\usepackage{xcolor}
\usepackage{hyperref}
\usepackage{geometry}
\usepackage{float}
\usepackage{microtype}

\geometry{margin=2.5cm}
\hypersetup{colorlinks=true,linkcolor=blue!70!black,citecolor=green!50!black,urlcolor=blue!70!black}

\title{\textbf{ReadSpyn}: A Semi-Classical Simulator for\\
  RF Reflectometry Readout of Quantum Dot Systems\\[4pt]
  \large Model Summary, Physical Assumptions, and\\
  Path Towards Realistic Simulation}
\author{Jan A.\ Krzywda \and Rouven K.\ Koch}
\date{\today}

\begin{document}
\maketitle

\begin{abstract}
We present \textsc{ReadSpyn}, a time-domain simulator for RF reflectometry
readout of semiconductor quantum dot (QD) systems.
The simulator propagates the equations of motion of an RLC resonator
capacitively coupled to one or more quantum dots, driven at its resonant
frequency, and subject to charge noise.
The reflected voltage is subsequently IQ-demodulated to yield the quadrature
signals $I(t)$ and $Q(t)$ recorded in experiment.
\textsc{ReadSpyn} represents a step towards realistically reproducing
laboratory readout traces: it incorporates an experimentally motivated
charge-noise spectrum, an analytic Coulomb-peak conductance model, and a
self-consistent circuit ODE.
We explicitly enumerate the physical approximations made and identify the
extensions required to close the gap with state-of-the-art experimental
conditions.
\end{abstract}

\tableofcontents
\newpage

%----------------------------------------------------------------------
\section{Introduction and Motivation}

Radio-frequency (RF) reflectometry is the workhorse readout technique for
spin and charge qubits in semiconductor QD devices
\cite{Schoelkopf1998,Ares2016,Gonzalez-Zalba2015}.
A resonant circuit is tank-matched to a sensor QD; changes in the sensor
conductance shift the resonator's impedance and are imprinted on the phase
and amplitude of the reflected microwave tone.
After IQ demodulation the two quadratures $I(t)$ and $Q(t)$ are used to
discriminate charge states in real time.

Simulating this measurement chain end-to-end is valuable because
(i) it allows optimisation of device parameters before fabrication,
(ii) it provides ground-truth data for benchmarking machine-learning
classifiers, and
(iii) it enables systematic studies of how different noise sources limit
readout fidelity.

\textsc{ReadSpyn} targets these use cases within a semi-classical framework
that is tractable and interpretable while capturing the dominant physics.
This document summarises the model, states the assumptions explicitly, and
outlines the road map towards higher-fidelity simulation.

%----------------------------------------------------------------------
\section{Quantum Dot System}

\subsection{Constant-Interaction Model}

The electrostatics of an array of $N$ quantum dots are described within the
constant-interaction (CI) model.
Mutual capacitances between dots and between dots and sensors are encoded in
the matrices $C_{dd} \in \mathbb{R}^{N\times N}$ and
$C_{ds} \in \mathbb{R}^{N\times M}$, where $M$ is the number of sensor QDs.

\subsection{Effective Detuning}

Given a charge-state vector $\mathbf{n}\in\mathbb{Z}^{N}$, a sensor-gate
voltage vector $\mathbf{V}_s\in\mathbb{R}^{M}$, and a global offset
$\varepsilon_0$, the effective detuning of sensor $j$ is
\begin{equation}
  \varepsilon_j
  = \bigl[C_{ds}^{\top}C_{dd}^{-1}\bigr]_{j}
    \bigl(\mathbf{n} + C_{ds}\mathbf{V}_s\bigr)
    + \varepsilon_0 \,.
  \label{eq:detuning}
\end{equation}
The row vector $[C_{ds}^{\top}C_{dd}^{-1}]_j$ is the \emph{lever arm}
projecting the charge configuration onto the detuning axis of sensor $j$.
Different charge states $\mathbf{n}^{(A)}$ and $\mathbf{n}^{(B)}$ produce
different detunings $\varepsilon^{(A)}_j$ and $\varepsilon^{(B)}_j$, which
is the physical origin of the IQ contrast used for charge-state
discrimination.

%----------------------------------------------------------------------
\section{Sensor Conductance: the Coulomb Peak}

The sensor QD is operated near a Coulomb blockade resonance.
In the sequential-tunnelling limit at low temperature the zero-bias
conductance as a function of detuning is \cite{Beenakker1991}
\begin{equation}
  G(\varepsilon) = \frac{2}{R_0}
    \cosh^{-2}\!\!\left(\frac{2\varepsilon}{\varepsilon_w}\right),
  \label{eq:coulomb_peak}
\end{equation}
where $R_0=1/g_0$ is the on-resonance resistance and $\varepsilon_w$ is the
thermal broadening width.
The local slope
\begin{equation}
  \frac{dG}{d\varepsilon}
  = -\frac{4}{R_0\,\varepsilon_w}
    \cosh^{-2}\!\!\left(\frac{2\varepsilon}{\varepsilon_w}\right)
    \tanh\!\left(\frac{2\varepsilon}{\varepsilon_w}\right)
  \label{eq:dGdeps}
\end{equation}
is zero at the peak centre and extremal at $|\varepsilon|=\varepsilon_w/2$.
A charge state whose static detuning \eqref{eq:detuning} lands on the steep
flank therefore produces larger IQ fluctuations under charge noise than one
sitting near the peak top --- even when both states experience the same noise.
This is the mechanism behind the asymmetric noise sensitivity observed in the
two-dot simulations.

%----------------------------------------------------------------------
\section{RLC Resonator and Circuit ODE}

\subsection{Lumped-Element Circuit}

The sensor QD is embedded in a lumped RLC tank circuit:
\begin{itemize}[nosep]
  \item Inductor $L_c$ (bond wire or on-chip spiral),
  \item Parasitic shunt capacitance $C_p$,
  \item Transmission-line load $R_L$ (typically $50\,\Omega$),
  \item Optional series coupling resistor $R_c$.
\end{itemize}
The circuit is driven at $\omega_d\approx\omega_0=1/\sqrt{L_cC_p}$.

\subsection{Equations of Motion}

With state vector $\mathbf{y}=(V_{C_p},I_L)^{\top}$ the circuit obeys
\begin{align}
  \frac{dV_{C_p}}{dt}
    &= \frac{I_L - G_s(t)V_{C_p}}{C_p}
       \qquad(R_c=0), \label{eq:ode1}\\[4pt]
  \frac{dI_L}{dt}
    &= \frac{V_s(t) - R_LI_L - V_A}{L_c}, \label{eq:ode2}
\end{align}
where $V_s(t)=V_0\cos(\omega_d t)$ and $V_A=V_{C_p}$ for $R_c=0$.
The time-dependent conductance $G_s(t)=G(\varepsilon(t))$ is obtained by
substituting the instantaneous noisy detuning into \eqref{eq:coulomb_peak}.
The ODEs are integrated with the stiff Radau solver
(\texttt{scipy.integrate.solve\_ivp}); inner-loop functions are JIT-compiled
with Numba.

\subsection{IQ Demodulation}

The voltage $V_A(t)$ is demodulated against the drive tone:
\begin{align}
  I(t) &= \mathrm{LPF}\!\left[V_A(t)\cos(\omega_d t+\phi)\right],\\
  Q(t) &= \mathrm{LPF}\!\left[V_A(t)\sin(\omega_d t+\phi)\right],
\end{align}
where $\phi$ is the demodulation phase.
In \textsc{ReadSpyn} demodulation uses the Hilbert transform; edge artefacts
are suppressed by padding the integration window before trimming to the
requested output interval.

%----------------------------------------------------------------------
\section{Noise Models}

\subsection{Spectrum-Prescribed $1/f$ Noise}

The dominant low-frequency noise in semiconductor QD devices is charge noise
with PSD $S(f)\propto 1/f$.
\textsc{ReadSpyn} generates noise trajectories $\delta\varepsilon(t)$ by
drawing Gaussian random amplitudes in frequency space scaled by $\sqrt{S(f)}$
and inverse-Fourier-transforming, then normalising to a prescribed RMS
amplitude $\sigma$ (in units of $\varepsilon_w$).
The same realisation can be shared across charge states (common-mode noise)
or independent realisations can be used.

\subsection{Ornstein--Uhlenbeck and Telegraph Noise}

For coloured Gaussian noise with a Lorentzian PSD, an Ornstein--Uhlenbeck
process is available:
\begin{equation}
  d\delta\varepsilon = -\gamma\,\delta\varepsilon\,dt
    + \sigma\sqrt{2\gamma}\,dW, \label{eq:ou}
\end{equation}
with correlation time $\tau_c=1/\gamma$.
Two-level fluctuators (bistable charge traps) are modelled as telegraph noise
with switching rate $\gamma$.

%----------------------------------------------------------------------
\section{A Step Towards Realistic Simulation}

\textsc{ReadSpyn} advances beyond a simple analytic toy model in several
respects:
\begin{enumerate}[label=\arabic*.]
  \item \textbf{Self-consistent circuit dynamics.} The conductance back-action
    on the resonator is solved via a full ODE rather than a perturbative
    susceptibility formula.
  \item \textbf{Time-domain noise injection.} Noise enters as a time-domain
    trajectory, so simulated IQ traces include realistic low-frequency drift.
  \item \textbf{User-configurable spectral shape.} The $1/f$ PSD can be matched
    to measured charge-noise spectra.
  \item \textbf{Multi-dot, multi-sensor geometry.} The capacitance-matrix
    formalism supports arbitrary numbers of dots and sensors, enabling studies
    of cross-talk and lever-arm asymmetries.
  \item \textbf{Edge-corrected Hilbert demodulation.} Window padding suppresses
    Hilbert-transform edge ringing, mirroring experimental DSP practice.
  \item \textbf{Stationary initial state.} The ODE is initialised at the
    ring-up-free steady state, avoiding transient artefacts.
\end{enumerate}

%----------------------------------------------------------------------
\section{Limitations with Respect to Realistic Semiconductor QD Measurements}
\label{sec:limitations}

Despite the advances above, several important physical effects are absent from
the current framework. We group them thematically below.

\subsection{Quantum and Many-Body Effects}

\begin{enumerate}[label=\alph*.]
  \item \textbf{No tunnel broadening.} Eq.~\eqref{eq:coulomb_peak} is the
    sequential-tunnelling result valid at $k_BT\gg\Gamma$. When $\Gamma\sim
    k_BT$ the Lorentzian (co-tunnelling / Kondo) line shape is not captured.
  \item \textbf{Classical charge states only.} Coherent charge superpositions
    (e.g.\ spin-to-charge conversion in Pauli-spin-blockade readout) require a
    density-matrix description.
  \item \textbf{No spin, valley, or orbital degeneracy.} Zeeman splitting,
    spin-orbit coupling, and valley degeneracy in Si/Ge are absent.
  \item \textbf{Constant-interaction approximation.} Energy-dependent
    screening and higher-order Coulomb effects are neglected.
\end{enumerate}

\subsection{Thermal and Non-Equilibrium Effects}

\begin{enumerate}[label=\alph*.]
  \item \textbf{Fixed electron temperature.} Microwave heating of the electron
    bath (photon-assisted tunnelling) is not modelled.
  \item \textbf{No photon-number-dependent back-action.} Above the
    single-photon power limit the resonator photon number causes an ac-Stark
    shift and measurement-induced dephasing, absent in the classical ODE.
  \item \textbf{No charge-state relaxation.} The charge configuration
    $\mathbf{n}$ is held constant during integration; qubit $T_1$ processes
    are not simulated.
\end{enumerate}

\subsection{Circuit and Electromagnetic Effects}

\begin{enumerate}[label=\alph*.]
  \item \textbf{Lumped-element model.} Distributed transmission-line effects
    and parasitic resonances of bond wires become important above
    ${\sim}1$\,GHz.
  \item \textbf{No amplifier noise.} The noise floor in experiment is
    dominated by the cryogenic amplifier (HEMT or paramp); this is not added
    to the simulated signal.
  \item \textbf{Perfect impedance matching.} Reflections from mismatch between
    the tank and the coaxial line are ignored.
  \item \textbf{No dispersive frequency shift.} Only the dissipative part
    $\delta G_s$ is modelled; the reactive part $\delta B_s$ that shifts the
    resonant frequency (dispersive readout) is absent.
\end{enumerate}

\subsection{Noise Modelling}

\begin{enumerate}[label=\alph*.]
  \item \textbf{Scalar, additive, global charge noise.} The noise is a single
    shift $\delta\varepsilon(t)$ on the sensor detuning. Per-gate voltage
    fluctuations that perturb individual dot potentials independently are not
    resolved.
  \item \textbf{No spatial cross-correlations.} Correlated noise from a
    single nearby charge trap affecting multiple dots is not included.
  \item \textbf{No TLS noise in circuit components.} Two-level-system noise
    in the dielectric of $C_p$ causes resistance fluctuations that are not
    modelled.
\end{enumerate}

\subsection{Signal Processing}

\begin{enumerate}[label=\alph*.]
  \item \textbf{Ideal local oscillator.} Phase noise and frequency jitter of
    the drive source are absent.
  \item \textbf{Ideal brick-wall low-pass filter.} The Hilbert-transform
    demodulation acts as a brick-wall filter; a realistic Butterworth or
    Gaussian filter would roll off differently.
  \item \textbf{No ADC effects.} Quantisation noise, dynamic range
    limitations, and clock jitter are not included.
\end{enumerate}

%----------------------------------------------------------------------
\section{Road Map}

Table~\ref{tab:roadmap} summarises the most impactful extensions and their
estimated implementation effort.

\begin{table}[H]
\centering
\renewcommand{\arraystretch}{1.3}
\begin{tabular}{lll}
\toprule
Extension & Physical motivation & Effort \\
\midrule
Amplifier noise (HEMT / paramp)           & Sets SNR floor in all experiments       & Low    \\
Finite electron temperature $T_e$          & Low-power / high-fidelity regime        & Low    \\
Per-dot local gate noise                   & Realistic multi-dot operation           & Low    \\
Charge relaxation / $T_1$ during readout   & Spin-qubit readout fidelity             & Medium \\
Dispersive frequency shift ($\delta B_s$)  & Dispersive readout schemes              & Medium \\
Tunnel-broadened Lorentzian line shape     & $\Gamma\sim k_BT$ regime               & Medium \\
Density-matrix charge dynamics             & Coherent charge superpositions          & High   \\
Distributed circuit / transmission line    & $f_0>1$\,GHz, parasitic resonances     & High   \\
\bottomrule
\end{tabular}
\caption{Prioritised road map for extending \textsc{ReadSpyn}.}
\label{tab:roadmap}
\end{table}

%----------------------------------------------------------------------
\section{Conclusion}

\textsc{ReadSpyn} provides a self-consistent, time-domain simulation of RF
reflectometry readout of semiconductor quantum dot systems.
By solving the circuit ODE in the presence of physically motivated charge
noise and demodulating the output, the simulator reproduces the qualitative
features of experimental IQ traces: state-dependent DC offsets, noise-driven
fluctuations whose magnitude depends on the local slope \eqref{eq:dGdeps} of
the Coulomb peak, and comet-tail trajectories in the IQ plane.

The framework is a deliberate step towards realistic measurement simulation
rather than a full first-principles model.
Section~\ref{sec:limitations} enumerates the current approximations in order
of physical importance.
The most pressing gaps are the absence of amplifier noise (which sets the
actual SNR floor) and of charge-relaxation dynamics (which limit spin-qubit
readout fidelity).
Addressing these two items, together with per-dot local gate noise for
multi-dot arrays, would bring \textsc{ReadSpyn} into quantitative contact
with state-of-the-art experimental data.

%----------------------------------------------------------------------
\begin{thebibliography}{9}

\bibitem{Schoelkopf1998}
R.~J.\ Schoelkopf \textit{et al.},
\textit{The radio-frequency single-electron transistor (RF-SET)},
Science \textbf{280}, 1238 (1998).

\bibitem{Ares2016}
N.\ Ares \textit{et al.},
\textit{Sensitive radio-frequency measurements of a quantum dot by tuning to
perfect impedance matching},
Phys.\ Rev.\ Applied \textbf{5}, 034011 (2016).

\bibitem{Gonzalez-Zalba2015}
M.~F.\ Gonzalez-Zalba \textit{et al.},
\textit{Probing the limits of gate-based charge sensing},
Nature Communications \textbf{6}, 6084 (2015).

\bibitem{Beenakker1991}
C.~W.~J.\ Beenakker,
\textit{Theory of Coulomb-blockade oscillations in the conductance of a
quantum dot},
Phys.\ Rev.\ B \textbf{44}, 1646 (1991).

\end{thebibliography}

\end{document}
